\hypertarget{ejercicio-2.-de-expresiuxf3n-regular-a-autuxf3mata}{%
\section{Ejercicio 2. De expresión regular a autómata}\label{ejercicio-2.-de-expresiuxf3n-regular-a-autuxf3mata}}

\hypertarget{notaciuxf3n-aceptada}{%
\subsection{Notación aceptada}\label{notaciuxf3n-aceptada}}

El analizador (\texttt{src/nucleo/expresiones.py}) acepta la notación del curso:

\begin{longtable}[]{@{}
  >{\raggedright\arraybackslash}p{(\columnwidth - 2\tabcolsep) * \real{0.5000}}
  >{\raggedright\arraybackslash}p{(\columnwidth - 2\tabcolsep) * \real{0.5000}}@{}}
\caption{Notación del analizador. Los espacios se ignoran.}\tabularnewline
\toprule\noalign{}
\begin{minipage}[b]{\linewidth}\raggedright
Escritura
\end{minipage} & \begin{minipage}[b]{\linewidth}\raggedright
Significado
\end{minipage} \\
\midrule\noalign{}
\endfirsthead
\toprule\noalign{}
\begin{minipage}[b]{\linewidth}\raggedright
Escritura
\end{minipage} & \begin{minipage}[b]{\linewidth}\raggedright
Significado
\end{minipage} \\
\midrule\noalign{}
\endhead
\bottomrule\noalign{}
\endlastfoot
\texttt{a}, \texttt{0}, \texttt{@}, \ldots{} & Cualquier carácter que no sea operador es un símbolo. \\
\texttt{\textbackslash{}x} & El carácter x literal aunque sea operador (por ejemplo \texttt{\textbackslash{}+}, \texttt{\textbackslash{}(}). \\
\texttt{λ} o \texttt{ε} & Palabra vacía. \\
\texttt{∅} o \texttt{Ø} & Lenguaje vacío. \\
\texttt{r\ +\ s} & Unión (precedencia más baja). \\
\texttt{rs} & Concatenación por yuxtaposición. \\
\texttt{r*} & Cerradura de Kleene (precedencia más alta). \\
\texttt{(\ r\ )} & Agrupación. \\
\texttt{{[}D{]}} \texttt{{[}m{]}} \texttt{{[}M{]}} \texttt{{[}L{]}} \texttt{{[}X{]}} & Clases declaradas: dígitos 0--9, minúsculas a--z, mayúsculas A--Z, letras y letras con dígitos. Se pueden registrar más con \texttt{Clases.registrar}. \\
\end{longtable}

Una clase es una abreviatura de la unión de sus símbolos: {[}D{]} = 0 + 1 + \ldots{} + 9.

\hypertarget{anuxe1lisis-shunting-yard-con-posiciones}{%
\subsection{\texorpdfstring{Análisis: \emph{shunting-yard} con posiciones}{Análisis: shunting-yard con posiciones}}\label{anuxe1lisis-shunting-yard-con-posiciones}}

El análisis tiene cuatro pasos, y cada ficha conserva la posición del carácter que la originó:

\begin{enumerate}
\def\labelenumi{\arabic{enumi}.}
\tightlist
\item
  \textbf{Léxico.} \texttt{Lexico} convierte el texto en fichas: símbolo, clase, λ, ∅, \texttt{+}, \texttt{*}, \texttt{(} y \texttt{)}.
\item
  \textbf{Concatenación implícita.} Entre una ficha que puede terminar un operando (símbolo, clase, λ, ∅, \texttt{)} o \texttt{*}) y una que puede empezarlo (símbolo, clase, λ, ∅ o \texttt{(}) se inserta un operador explícito \texttt{·}. En este mismo recorrido se detectan los errores de forma: un \texttt{+} sin operando, una \texttt{*} que no se aplica a nada o paréntesis vacíos.
\item
  \textbf{Posfija.} El algoritmo \emph{shunting-yard} de Dijkstra pasa la secuencia a notación posfija con una pila de operadores, respetando \texttt{*} \textgreater{} \texttt{·} \textgreater{} \texttt{+} y los paréntesis; aquí se detectan los paréntesis desbalanceados.
\item
  \textbf{Árbol.} La posfija se reduce con una pila de nodos: \texttt{Simbolo}, \texttt{Clase}, \texttt{Lambda}, \texttt{Vacio}, \texttt{Union}, \texttt{Concatenacion} y \texttt{Estrella}.
\end{enumerate}

Cada nodo sabe escribirse en la notación del curso con los paréntesis mínimos (\texttt{notacion()}), traducirse a la sintaxis de \texttt{re} (\texttt{python()}), decir si es anulable y qué símbolos usa.

\begin{longtable}[]{@{}ll@{}}
\caption{Expresiones mal formadas y su mensaje.}\tabularnewline
\toprule\noalign{}
Entrada & Mensaje \\
\midrule\noalign{}
\endfirsthead
\toprule\noalign{}
Entrada & Mensaje \\
\midrule\noalign{}
\endhead
\bottomrule\noalign{}
\endlastfoot
\texttt{(a\ +\ b} & Falta el «)» que cierra este paréntesis (posición 1) \\
\texttt{a\ +\ *\ b} & La «*» no se aplica a nada (posición 5) \\
\texttt{ab)} & «)» sin un «(» que lo abra (posición 3) \\
\texttt{a()} & Paréntesis vacíos; use λ para la palabra vacía (posición 2) \\
\texttt{{[}Z{]}a} & La clase «{[}Z{]}» no está declarada (posición 1) \\
\texttt{a\ +} & Al «+» le falta el operando derecho (posición 3) \\
\end{longtable}

En la interfaz la revisión ocurre \textbf{mientras se escribe}: el campo muestra el mensaje y debajo aparece la expresión con un \texttt{\^{}} bajo el carácter culpable.

\evidencia[0.95]{../evidencias/practica4/app/01-error-posicion.png}{ER → AFD: error de sintaxis señalado con su posición.}

\hypertarget{parte-i-del-teorema-de-kleene-paso-a-paso}{%
\subsection{Parte I del teorema de Kleene, paso a paso}\label{parte-i-del-teorema-de-kleene-paso-a-paso}}

\texttt{ConstruccionThompson} recorre el árbol en posorden y produce un \textbf{fragmento} (estado inicial, estado final, estados y aristas) por nodo. Los estados se numeran q0, q1, \ldots{} en el orden en que se crean, de modo que un estado conserva su nombre en todas las etapas.

\begin{longtable}[]{@{}
  >{\raggedright\arraybackslash}p{(\columnwidth - 2\tabcolsep) * \real{0.2000}}
  >{\raggedright\arraybackslash}p{(\columnwidth - 2\tabcolsep) * \real{0.8000}}@{}}
\caption{Casos básicos y construcciones \autocites[sec.~1.3, lema 1.55]{sipser2013}[sec.~3.2.3]{hopcroft2007}.}\tabularnewline
\toprule\noalign{}
\begin{minipage}[b]{\linewidth}\raggedright
Nodo
\end{minipage} & \begin{minipage}[b]{\linewidth}\raggedright
Construcción
\end{minipage} \\
\midrule\noalign{}
\endfirsthead
\toprule\noalign{}
\begin{minipage}[b]{\linewidth}\raggedright
Nodo
\end{minipage} & \begin{minipage}[b]{\linewidth}\raggedright
Construcción
\end{minipage} \\
\midrule\noalign{}
\endhead
\bottomrule\noalign{}
\endlastfoot
∅ & Dos estados nuevos i, f sin transiciones. \\
λ & i →λ f. \\
a & i →a f. \\
Clase & i → f con una transición por cada símbolo de la clase. \\
r + s & Nuevos i, f; i →λ inicio(r), i →λ inicio(s); fin(r) →λ f, fin(s) →λ f. \\
rs & fin(r) →λ inicio(s); el inicial es el de r y el final el de s. \\
r* & Nuevos i, f; i →λ inicio(r), i →λ f, fin(r) →λ inicio(r), fin(r) →λ f. \\
\end{longtable}

Cada construcción queda registrada como una etapa con su autómata parcial; la sección \textbf{ER → AFD} las recorre con «Anterior / Siguiente», dibuja cada fragmento y escribe qué pieza se está construyendo («Unión: a + b», «Estrella: (a + b)*»). Para (ab)*a (272b) las seis etapas son: caso básico a; caso básico b; concatenación ab; estrella (ab)*; caso básico a; concatenación (ab)*a, con 8 estados en total.

A ese AFN-λ se le aplican los algoritmos de la Práctica 3 ---eliminación de λ, subconjuntos y minimización--- y la misma barra de etapas permite ver cada uno. Cada autómata (AFN-λ, AFN, AFD y AFD mínimo) se exporta a \texttt{.jff} con \emph{Guardar .jff} o los cuatro a la vez con \emph{Guardar las cuatro}.

\evidencia[0.95]{../evidencias/practica4/app/02-thompson.png}{ER → AFD con (a + b)\emph{abb: etapa «Estrella: (a + b)}» de la
construcción.}

\evidencia[0.95]{../evidencias/practica4/app/03-minimo.png}{El AFD mínimo de (a + b)*abb y la prueba de la palabra \texttt{babb}.}

\hypertarget{validaciuxf3n-de-palabras}{%
\subsection{Validación de palabras}\label{validaciuxf3n-de-palabras}}

La tarjeta \emph{Probar palabras} valida mientras se escribe con el AFD mínimo y muestra también el resultado de \texttt{re.fullmatch} con la traducción de la expresión; si alguna vez difirieran, aparecería una advertencia.

\hypertarget{casos-de-prueba}{%
\subsection{Casos de prueba}\label{casos-de-prueba}}

\begin{longtable}[]{@{}
  >{\raggedright\arraybackslash}p{(\columnwidth - 10\tabcolsep) * \real{0.1667}}
  >{\raggedright\arraybackslash}p{(\columnwidth - 10\tabcolsep) * \real{0.1667}}
  >{\centering\arraybackslash}p{(\columnwidth - 10\tabcolsep) * \real{0.1667}}
  >{\centering\arraybackslash}p{(\columnwidth - 10\tabcolsep) * \real{0.1667}}
  >{\centering\arraybackslash}p{(\columnwidth - 10\tabcolsep) * \real{0.1667}}
  >{\centering\arraybackslash}p{(\columnwidth - 10\tabcolsep) * \real{0.1667}}@{}}
\caption{Tamaño de los autómatas en cada etapa (los AFD no incluyen el subconjunto ∅).}\tabularnewline
\toprule\noalign{}
\begin{minipage}[b]{\linewidth}\raggedright
Caso
\end{minipage} & \begin{minipage}[b]{\linewidth}\raggedright
Expresión
\end{minipage} & \begin{minipage}[b]{\linewidth}\centering
Construcciones
\end{minipage} & \begin{minipage}[b]{\linewidth}\centering
AFN-λ
\end{minipage} & \begin{minipage}[b]{\linewidth}\centering
AFD
\end{minipage} & \begin{minipage}[b]{\linewidth}\centering
AFD mínimo
\end{minipage} \\
\midrule\noalign{}
\endfirsthead
\toprule\noalign{}
\begin{minipage}[b]{\linewidth}\raggedright
Caso
\end{minipage} & \begin{minipage}[b]{\linewidth}\raggedright
Expresión
\end{minipage} & \begin{minipage}[b]{\linewidth}\centering
Construcciones
\end{minipage} & \begin{minipage}[b]{\linewidth}\centering
AFN-λ
\end{minipage} & \begin{minipage}[b]{\linewidth}\centering
AFD
\end{minipage} & \begin{minipage}[b]{\linewidth}\centering
AFD mínimo
\end{minipage} \\
\midrule\noalign{}
\endhead
\bottomrule\noalign{}
\endlastfoot
L3-241 & 0 + 1 & 3 & 6 & 3 & 2 \\
L3-242 & 0* + 1 & 4 & 8 & 3 & 3 \\
L3-243 & (0 + 1)* & 4 & 8 & 3 & 1 \\
L3-244 & 0(0 + 1) & 5 & 8 & 4 & 3 \\
L3-245 & 0(0 + 1)* & 6 & 10 & 4 & 2 \\
L3-246 & 0*(0 + 1) & 6 & 10 & 3 & 3 \\
L3-247 & (01)* + 1* & 7 & 12 & 4 & 4 \\
L3-248 & (0* + 1)0 & 6 & 10 & 4 & 4 \\
L3-249 & (0* + 1)0* & 7 & 12 & 4 & 2 \\
L3-250 & 0*1* + 1*0 & 10 & 16 & 6 & 5 \\
L3-251 & (0 + 1*)0 + 1*1 & 11 & 18 & 4 & 4 \\
L3-252 & (0 + 1)*(00 + 11) + (00 + 11)(0 + 1)* & 25 & 38 & 13 & 8 \\
L3-253 & (0 + (11 + 00)*)(11 + 10 + 01)*(λ + 110) & 31 & 44 & 22 & 17 \\
L4-272a & a + b* & 4 & 8 & 3 & 3 \\
L4-272b & (ab)*a & 6 & 8 & 3 & 2 \\
L4-272c & (a + b)*abb & 10 & 14 & 5 & 4 \\
L4-272d & a(ba)* + b & 8 & 12 & 5 & 4 \\
L4-272e & (a + λ)(b + ∅) & 7 & 12 & 3 & 3 \\
L4-278 & (a + b)*b & 6 & 10 & 3 & 2 \\
\end{longtable}

\textbf{Lista 4, 273.} Procesando por λ-clausuras el AFN-λ de (ab)*a se aceptan \texttt{aba} y se rechazan \texttt{ab} y \texttt{abab}; con el de (a + b)*abb se acepta \texttt{babb} y se rechazan \texttt{aba}, \texttt{ab} y \texttt{abab}. Por ejemplo, en el segundo la lectura de \texttt{babb} empieza en λ-cl(q6) = \{q0, q2, q4, q6, q7, q8\} y termina en un conjunto que contiene al final q13.

\textbf{Lista 4, 278.} Para (a + b)*b: el AFN-λ tiene 10 estados; al quitar λ queda un AFN de 10 estados; la construcción de subconjuntos da 3 estados alcanzables y la minimización 2 (uno para «la última letra fue b», final, y otro para el resto). La eliminación de estados sobre ese AFD mínimo da a*b(b + aa*b)* y el lema de Arden (a + bb*a)*bb*; las dos son equivalentes a (a + b)*b según la prueba de equivalencia, con lo que se cierra el ciclo.

\textbf{Pruebas.} \texttt{tests/test\_expresiones.py} (81 pruebas) comprueba la precedencia, la notación con paréntesis mínimos (y que es estable al volver a analizarla), quince expresiones mal formadas con su posición exacta, la traducción a \texttt{re}, la forma de los autómatas de Thompson (un final sin salidas, a lo sumo dos estados por nodo) y, para cada uno de los 19 casos, que \textbf{el AFN-λ, el AFN, el AFD y el AFD mínimo coinciden en todas las palabras de longitud ≤ 10} y que el AFD mínimo coincide con \texttt{re.fullmatch} en todas las de longitud ≤ 7.
